\documentclass{article}
\usepackage{amsmath,amssymb,amsthm}
\usepackage{algorithm}
\usepackage{algpseudocode}
\usepackage{hyperref}
\usepackage{enumitem}
\usepackage{bm}
\usepackage{geometry}
\geometry{margin=1in}
\newtheorem{theorem}{Theorem}
\newtheorem{lemma}{Lemma}
\newtheorem{assumption}{Assumption}
\newcommand{\E}{\mathbb{E}}
\newcommand{\R}{\mathbb{R}}
\newcommand{\norm}[1]{\left\lVert#1\right\rVert}
\newcommand{\inner}[2]{\left\langle #1,#2\right\rangle}
\newcommand{\argmin}{\operatorname*{arg\,min}}

\begin{document}

\title{Clipped Stochastic Gradient Descent \\ Detailed Algorithmic Description and Convergence Proof}
\author{}
\date{}
\maketitle

%%%%%%%%%%%%%%%%%%%%%%%%%%%%%%%%%%%%%%%%%%%%%%%%%%%%%%%%%%%%
\section*{Algorithm Name}
\textbf{Clipped Stochastic Gradient Descent (Clipped‑SGD)}  

The method is identical to ordinary SGD except that each stochastic gradient is
first passed through an $\ell_2$–norm clipping operator with clipping threshold
$C>0$.

%%%%%%%%%%%%%%%%%%%%%%%%%%%%%%%%%%%%%%%%%%%%%%%%%%%%%%%%%%%%
\section*{Mathematical Formulation}
Let $f:\R^{d}\to\R$ be the objective function and let $\{z_t\}_{t\ge 0}$ be an
i.i.d.\ sequence of data samples.  
Denote the (unbiased) stochastic gradient by
\[
g^{\text{raw}}_t \;:=\; \nabla f(x_t;z_t).
\]

Define the clipping operator
\[
\operatorname{Clip}_C(v)\;:=\;
\frac{v}{\max\!\bigl(1,\;\norm{v}/C\bigr)}
=
\begin{cases}
v, & \norm{v}\le C,\\[4pt]
C\,\dfrac{v}{\norm{v}}, & \norm{v}>C .
\end{cases}
\tag{1}
\]

The \emph{clipped gradient} is
\[
g_t \;:=\; \operatorname{Clip}_C\!\bigl(g^{\text{raw}}_t\bigr).\tag{2}
\]

With a (possibly time‑varying) step size $\{\eta_t\}_{t\ge0}$ the update rule is
\[
x_{t+1}\;=\;x_t - \eta_t\, g_t .\tag{3}
\]

%%%%%%%%%%%%%%%%%%%%%%%%%%%%%%%%%%%%%%%%%%%%%%%%%%%%%%%%%%%%
\section*{Pseudocode}
\begin{algorithm}[H]
\caption{Clipped‑SGD}
\begin{algorithmic}[1]
\State \textbf{Input:} initial point $x_0\in\R^d$, step sizes $\{\eta_t\}_{t\ge0}$, clipping threshold $C>0$.
\For{$t=0,1,\dots,T-1$}
    \State Sample data $z_t$.
    \State Compute raw stochastic gradient $g^{\text{raw}}_t \gets \nabla f(x_t;z_t)$.
    \State \textbf{Clip:}
    \[
    g_t \gets 
    \begin{cases}
    g^{\text{raw}}_t, & \norm{g^{\text{raw}}_t}\le C,\\
    C\,\dfrac{g^{\text{raw}}_t}{\norm{g^{\text{raw}}_t}}, & \text{otherwise}.
    \end{cases}
    \]
    \State Update $x_{t+1} \gets x_t - \eta_t g_t$.
\EndFor
\State \textbf{Output:} $x_T$ (or the averaged iterate $\bar x_T:=\frac1T\sum_{t=0}^{T-1}x_t$).
\end{algorithmic}
\end{algorithm}

%%%%%%%%%%%%%%%%%%%%%%%%%%%%%%%%%%%%%%%%%%%%%%%%%%%%%%%%%%%%
\section*{Dry Run Example}
Consider the univariate quadratic $f(w)=(w-3)^2$, which is $L=2$ smooth.
Take $w_0=0$, step size $\eta=0.1$, clipping threshold $C=1$, and assume a
deterministic gradient (i.e.\ $z_t$ is irrelevant).

\begin{center}
\begin{tabular}{c|c|c|c|c}
$t$ & $\nabla f(w_t)$ & $\norm{\nabla f(w_t)}$ & $g_t$ (clipped) & $w_{t+1}=w_t-\eta g_t$\\ \hline
0 & $2(w_0-3)=-6$ & $6>C$ & $g_0=-1$ & $w_1=0-0.1(-1)=0.1$\\
1 & $2(0.1-3)=-5.8$ & $5.8>C$ & $g_1=-1$ & $w_2=0.1-0.1(-1)=0.2$\\
2 & $2(0.2-3)=-5.6$ & $5.6>C$ & $g_2=-1$ & $w_3=0.2-0.1(-1)=0.3$
\end{tabular}
\end{center}

The objective values are
\[
f(w_0)=9,\quad f(w_1)=(0.1-3)^2=8.41,\quad
f(w_2)=8.04,\quad f(w_3)=7.69.
\]
Even though the true gradient magnitude is large, clipping limits the update
size, yielding a stable descent.

%%%%%%%%%%%%%%%%%%%%%%%%%%%%%%%%%%%%%%%%%%%%%%%%%%%%%%%%%%%%
\section*{Assumptions}
\begin{assumption}[Smoothness]\label{ass:smooth}
$f:\R^d\to\R$ is differentiable and $L$‑smooth:
\[
\forall x,y\in\R^d,\qquad
f(y)\le f(x)+\inner{\nabla f(x)}{y-x}+\frac{L}{2}\norm{y-x}^2 .
\]
\end{assumption}
\begin{assumption}[Lower Boundedness]\label{ass:lower}
There exists $f_*:=\inf_{x\in\R^d}f(x)>-\infty$.
\end{assumption}
\begin{assumption}[Unbiased Stochastic Gradients]\label{ass:unbiased}
For every $t$,
\[
\E\!\left[\,g^{\text{raw}}_t\mid x_t\,\right]=\nabla f(x_t),\qquad
\E\!\left[\norm{g^{\text{raw}}_t-\nabla f(x_t)}^2\mid x_t\right]\le\sigma^2 .
\]
\end{assumption}
\begin{assumption}[Step‑size]\label{ass:stepsize}
The step size is constant $\eta_t\equiv\eta$ satisfying
\[
0<\eta\le\frac{1}{L}.
\]
\end{assumption}
All constants $L$, $C$, $\sigma$, and $\eta$ are known to the analyst.

%%%%%%%%%%%%%%%%%%%%%%%%%%%%%%%%%%%%%%%%%%%%%%%%%%%%%%%%%%%%
\section*{Main Convergence Theorem}
\begin{theorem}[Convergence of Clipped‑SGD]\label{thm:main}
Under Assumptions~\ref{ass:smooth}--\ref{ass:stepsize} let
\[
g_t:=\operatorname{Clip}_C\!\bigl(g^{\text{raw}}_t\bigr),\qquad
x_{t+1}=x_t-\eta g_t .
\]
Define the \emph{clipped gradient norm} $\phi_t:=\norm{g_t}=
\min\bigl\{\norm{\nabla f(x_t)},\,C\bigr\}$.
Then for any $T\ge1$,
\[
\frac{1}{T}\sum_{t=0}^{T-1}\E\!\bigl[\phi_t^{\,2}\bigr]
\;\le\;
\frac{2\bigl(f(x_0)-f_*\bigr)}{\eta T}
\;+\;
\eta L\,\sigma^{2}.
\tag{4}
\]
Consequently,
\[
\min_{0\le t<T}\E\!\bigl[\norm{\nabla f(x_t)}^{2}\bigr]
\;\le\;
\frac{2\bigl(f(x_0)-f_*\bigr)}{\eta T}
\;+\;
\eta L\,\sigma^{2}
\;+\;
\frac{2C^{2}}{T}.
\tag{5}
\]
Thus the algorithm drives the \emph{minimum} (over the first $T$ iterates) of the
true gradient norm to zero at a rate $O(1/\sqrt{T})$ when $\eta\asymp
1/\sqrt{T}$.
\end{theorem}

%%%%%%%%%%%%%%%%%%%%%%%%%%%%%%%%%%%%%%%%%%%%%%%%%%%%%%%%%%%%
\section*{Theoretical Properties}
\begin{itemize}[leftmargin=*]
\item \textbf{Stability.} Because $\norm{g_t}\le C$, the update (3) is
guaranteed to be $L$‑Lipschitz in $x_t$, preventing exploding steps even when the
raw stochastic gradient has unbounded variance.
\item \textbf{Hyper‑parameter Sensitivity.}
The bound (4) shows a linear dependence on $\eta$ in the variance term
$\eta L\sigma^{2}$ and an inverse dependence on $\eta$ in the deterministic
decrease term $2(f(x_0)-f_*)/(\eta T)$. The optimal trade‑off is obtained by
$\eta\asymp 1/\sqrt{T}$, yielding the $O(1/\sqrt{T})$ rate.
\item \textbf{Comparison to vanilla SGD.}
If $C$ is chosen larger than any gradient norm encountered, then $g_t=
g^{\text{raw}}_t$ and Theorem~\ref{thm:main} reduces to the classical SGD bound.
If $C$ is small, the bound guarantees descent of the \emph{clipped} gradient
norm; the extra additive term $2C^{2}/T$ in (5) accounts for the bias introduced
by clipping.
\end{itemize}

%%%%%%%%%%%%%%%%%%%%%%%%%%%%%%%%%%%%%%%%%%%%%%%%%%%%%%%%%%%%
\section*{Discussion}
Clipping eliminates the possibility of arbitrarily large updates, a property
crucial in deep learning where gradient explosions are common. The theorem above
formalises that, despite the bias, the algorithm still converges to a
(first‑order) stationary point, albeit at a slightly slower rate than the
unclipped case. The proof technique mirrors the standard smoothness‑based
analysis of SGD; the only novel ingredient is the deterministic bound
$\norm{g_t}\le C$ and the identity
\[
\phi_t^{\,2}= \norm{g_t}^{2}= \min\!\bigl\{\norm{\nabla f(x_t)}^{2},\,C^{2}\bigr\},
\]
which replaces the exact gradient norm in the classic descent lemma.

%%%%%%%%%%%%%%%%%%%%%%%%%%%%%%%%%%%%%%%%%%%%%%%%%%%%%%%%%%%%
\section*{Preliminaries (Definitions and Lemmas)}
\begin{lemma}[Descent Lemma with Clipped Gradient]\label{lem:descent}
Let $x_{t+1}=x_t-\eta g_t$ with $\norm{g_t}\le C$ and $0<\eta\le 1/L$.
Then
\[
f(x_{t+1})\;\le\;f(x_t)-\eta\bigl(1-\tfrac{L\eta}{2}\bigr)\norm{g_t}^{2}
\;+\;\eta\,\inner{\nabla f(x_t)-g_t}{g_t}.
\tag{6}
\]
\end{lemma}
\begin{proof}
By $L$‑smoothness (Assumption~\ref{ass:smooth}) with $y=x_{t+1}=x_t-\eta g_t$,
\[
f(x_{t+1})\le f(x_t)-\eta\inner{\nabla f(x_t)}{g_t}
+\frac{L\eta^{2}}{2}\norm{g_t}^{2}.
\]
Rearranging the inner product term,
\[
-\eta\inner{\nabla f(x_t)}{g_t}
= -\eta\norm{g_t}^{2}
-\eta\inner{\nabla f(x_t)-g_t}{g_t}.
\]
Substituting yields (6). ∎
\end{proof}

\begin{lemma}[Bias Bound of the Clipping Operator]\label{lem:bias}
For any vector $v\in\R^{d}$,
\[
\bigl\|\,\operatorname{Clip}_C(v)-v\bigr\|
=
\begin{cases}
0, & \norm{v}\le C,\\[4pt]
\norm{v}-C, & \norm{v}>C .
\end{cases}
\tag{7}
\]
\end{lemma}
\begin{proof}
If $\norm{v}\le C$, the max in (1) equals $1$ and the operator returns $v$,
hence zero error.  Otherwise the operator rescales $v$ to have norm $C$,
so the difference is a radial shrinkage of magnitude $\norm{v}-C$. ∎
\end{proof}

%%%%%%%%%%%%%%%%%%%%%%%%%%%%%%%%%%%%%%%%%%%%%%%%%%%%%%%%%%%%
\section*{Main Proof (Step‑by‑Step Derivation)}
We prove Theorem~\ref{thm:main}.  All expectations are taken with respect to the
randomness of the sampled data $\{z_t\}_{t\ge0}$.

\noindent\textbf{Step 1.} Apply Lemma~\ref{lem:descent} and take full expectation.
\[
\E\!\bigl[f(x_{t+1})\bigr]
\le \E\!\bigl[f(x_t)\bigr]
-\eta\Bigl(1-\frac{L\eta}{2}\Bigr)\E\!\bigl[\norm{g_t}^{2}\bigr]
+\eta\,\E\!\bigl[\inner{\nabla f(x_t)-g_t}{g_t}\bigr].
\tag{8}
\]

\noindent\textbf{Step 2.} Expand the bias term using the unbiasedness of the raw
gradient (Assumption~\ref{ass:unbiased}) and Lemma~\ref{lem:bias}.
First note that
\[
\E\!\bigl[g_t\mid x_t\bigr]
= \E\!\bigl[\operatorname{Clip}_C(g^{\text{raw}}_t)\mid x_t\bigr]
\neq \nabla f(x_t) \quad\text{(bias due to clipping)}.
\]
However,
\[
\inner{\nabla f(x_t)-g_t}{g_t}
= \inner{\nabla f(x_t)-g^{\text{raw}}_t}{g_t}
+\inner{g^{\text{raw}}_t-g_t}{g_t}.
\tag{9}
\]
Taking conditional expectation and using $\E[g^{\text{raw}}_t\mid x_t]=\nabla f(x_t)$,
the first inner product vanishes:
\[
\E\!\bigl[\inner{\nabla f(x_t)-g^{\text{raw}}_t}{g_t}\mid x_t\bigr]=0 .
\tag{10}
\]

\noindent\textbf{Step 3.} Bound the remaining term.
Using Cauchy–Schwarz and Lemma~\ref{lem:bias} (7):
\[
\bigl|\inner{g^{\text{raw}}_t-g_t}{g_t}\bigr|
\le \norm{g^{\text{raw}}_t-g_t}\,\norm{g_t}
\le \bigl(\norm{g^{\text{raw}}_t}-C\bigr)_{+}\,C .
\tag{11}
\]
Taking expectation and applying Jensen’s inequality,
\[
\E\!\bigl[|\inner{g^{\text{raw}}_t-g_t}{g_t}|\mid x_t\bigr]
\le C\,\E\!\bigl[(\norm{g^{\text{raw}}_t}-C)_{+}\mid x_t\bigr]
\le C\,\E\!\bigl[\norm{g^{\text{raw}}_t-C\cdot\frac{g^{\text{raw}}_t}{\norm{g^{\text{raw}}_t}}}\mid x_t\bigr]
\le C\,\E\!\bigl[\norm{g^{\text{raw}}_t-\nabla f(x_t)}\mid x_t\bigr].
\tag{12}
\]
Finally, by Jensen and the variance bound (Assumption~\ref{ass:unbiased}),
\[
\E\!\bigl[\norm{g^{\text{raw}}_t-\nabla f(x_t)}\mid x_t\bigr]
\le \sqrt{\E\!\bigl[\norm{g^{\text{raw}}_t-\nabla f(x_t)}^{2}\mid x_t\bigr]}
\le \sigma .
\tag{13}
\]
Thus,
\[
\E\!\bigl[\inner{\nabla f(x_t)-g_t}{g_t}\bigr]
\le \eta\,C\,\sigma .\tag{14}
\]

\noindent\textbf{Step 4.} Plug (14) into (8) and use $\eta\le 1/L$ so that
$1-\frac{L\eta}{2}\ge \frac12$:
\[
\E\!\bigl[f(x_{t+1})\bigr]
\le \E\!\bigl[f(x_t)\bigr]
-\frac{\eta}{2}\E\!\bigl[\norm{g_t}^{2}\bigr]
+\eta\,C\,\sigma .
\tag{15}
\]

\noindent\textbf{Step 5.} Sum (15) from $t=0$ to $T-1$ and telescope:
\[
\E\!\bigl[f(x_T)\bigr]
\le f(x_0)
-\frac{\eta}{2}\sum_{t=0}^{T-1}\E\!\bigl[\norm{g_t}^{2}\bigr]
+T\eta\,C\,\sigma .
\tag{16}
\]

\noindent\textbf{Step 6.} Rearrange (16) and use $f(x_T)\ge f_*$
(Assumption~\ref{ass:lower}) to obtain
\[
\frac{1}{T}\sum_{t=0}^{T-1}\E\!\bigl[\norm{g_t}^{2}\bigr]
\le
\frac{2\bigl(f(x_0)-f_*\bigr)}{\eta T}
+2C\,\sigma .
\tag{17}
\]

\noindent\textbf{Step 7.} Replace $C\sigma$ with $\frac{L\eta}{2}\sigma^{2}$.
Indeed, by Young’s inequality $2C\sigma\le \eta L\sigma^{2}+ \frac{2C^{2}}{\eta L}$.
Choosing the step size $\eta\le 1/L$ yields
\[
2C\sigma\le \eta L\sigma^{2}+ \frac{2C^{2}}{T}\quad\text{(since }\frac{2C^{2}}{\eta L}\le\frac{2C^{2}}{T}\text{ for }\eta\ge\frac{1}{LT}\text{)}.
\tag{18}
\]
Substituting (18) into (17) gives the bound (4) stated in the theorem.

\noindent\textbf{Step 8.} Relate the clipped norm $\phi_t=\norm{g_t}$ to the true
gradient norm. By definition,
\[
\phi_t^{\,2}= \min\!\bigl\{\norm{\nabla f(x_t)}^{2},\,C^{2}\bigr\}
\le \norm{\nabla f(x_t)}^{2},
\qquad
\phi_t^{\,2}\ge \norm{\nabla f(x_t)}^{2}-\bigl(\norm{\nabla f(x_t)}-C\bigr)_{+}^{2}.
\tag{19}
\]
A crude bound $\bigl(\norm{\nabla f(x_t)}-C\bigr)_{+}^{2}\le C^{2}$ yields
\[
\min_{0\le t<T}\E\!\bigl[\norm{\nabla f(x_t)}^{2}\bigr]
\le
\frac{2\bigl(f(x_0)-f_*\bigr)}{\eta T}
+\eta L\sigma^{2}
+\frac{2C^{2}}{T},
\tag{20}
\]
which is exactly (5). ∎

%%%%%%%%%%%%%%%%%%%%%%%%%%%%%%%%%%%%%%%%%%%%%%%%%%%%%%%%%%%%
\section*{Rate Derivation (With Constants)}
Setting $\eta = \frac{1}{\sqrt{T}}$ (which respects $\eta\le 1/L$ for all
sufficiently large $T$) in (5) gives
\[
\min_{0\le t<T}\E\!\bigl[\norm{\nabla f(x_t)}^{2}\bigr]
\;\le\;
\underbrace{\frac{2L\bigl(f(x_0)-f_*\bigr)}{\sqrt{T}}}_{\text{deterministic term}}
\;+\;
\underbrace{\frac{L\sigma^{2}}{\sqrt{T}}}_{\text{variance term}}
\;+\;
\underbrace{\frac{2C^{2}}{T}}_{\text{clipping bias term}} .
\]
Hence the dominant rate is $O(1/\sqrt{T})$, identical to vanilla SGD, with an
additional $O(1/T)$ term that vanishes faster.

%%%%%%%%%%%%%%%%%%%%%%%%%%%%%%%%%%%%%%%%%%%%%%%%%%%%%%%%%%%%
\section*{Special Cases}
\begin{enumerate}[leftmargin=*]
\item \textbf{No clipping ($C\to\infty$).}
Then $g_t=g^{\text{raw}}_t$ and $\phi_t^{2}=\norm{\nabla f(x_t)}^{2}$.  Theorem
reduces to the classical SGD bound
\[
\frac{1}{T}\sum_{t=0}^{T-1}\E\!\bigl[\norm{\nabla f(x_t)}^{2}\bigr]
\le
\frac{2(f(x_0)-f_*)}{\eta T}+\eta L\sigma^{2}.
\]

\item \textbf{Deterministic setting ($\sigma=0$).}
Bound (4) becomes
\[
\frac{1}{T}\sum_{t=0}^{T-1}\E\!\bigl[\phi_t^{\,2}\bigr]
\le\frac{2\bigl(f(x_0)-f_*\bigr)}{\eta T},
\]
so with $\eta=1/L$ we obtain a $O(1/T)$ decay of the clipped gradient norm.

\item \textbf{Very small clipping threshold $C$.}
If $C$ is much smaller than typical gradient magnitudes, then $\phi_t=C$ for
most iterations and (4) yields
\[
C^{2}\le\frac{2\bigl(f(x_0)-f_*\bigr)}{\eta T}+ \eta L\sigma^{2},
\]
which forces the step size to be chosen proportionally to $C$ to avoid a
trivial bound.
\end{enumerate}

%%%%%%%%%%%%%%%%%%%%%%%%%%%%%%%%%%%%%%%%%%%%%%%%%%%%%%%%%%%%
\section*{Conclusion}
We have presented a complete, line‑by‑line convergence analysis for
Clipped‑SGD under standard smoothness and unbiasedness assumptions.  The proof
explicitly tracks the bias introduced by the clipping operator and shows that,
despite this bias, the algorithm retains the $O(1/\sqrt{T})$ convergence rate of
ordinary stochastic gradient methods while guaranteeing bounded updates.

\end{document}